\chapter{从精确 Mie 到 Rayleigh：长波近似究竟近似了什么}
\label{chap:rayleigh}

Rayleigh 公式不另起一套静电理论，而应当从第~\ref{chap:mie-sphere} 章的精确 Mie 系数在尺寸参数 $x=k_{\rm h}a$ 的小量极限中展开。这样才能看清材料对比、几何尺寸和电子相位匹配分别来自哪里，也能明确 $x\sim1$ 时究竟是哪一项先失效。本章只在球形主线上完成这一极限；一般椭球和无限圆柱作为几何扩展留到附录。

\section{长波展开的统一尺度}

球形粒子的完整多极解确定以后，再施加长波条件，把精确 Maxwell 散射按尺寸参数 $x=k_{\rm h}a$ 展开。正文只沿单球主线推进：先从精确部分波阶数证明偶极为何占主导，再独立求球的准静态极化率，最后把两者逐项核对。一般椭球只作为附录中的几何推广，不参与球形 Rayleigh 结果的建立。电子端始终不需要另立新理论，只在最后把所得散射场代入式~\eqref{eq:Fparallel-general-def} 的同一轨迹投影。设粒子嵌在均匀、各向同性宿主中，粒子和宿主的相对介电函数、相对磁导率分别为
\begin{equation}
\varepsilon_{\rm p}(\omega),\quad \mu_{\rm p}(\omega),
\qquad
\varepsilon_{\rm h}(\omega),\quad \mu_{\rm h}(\omega).
\end{equation}
定义
\begin{equation}
\varepsilon_{\rm rel}\equiv\frac{\varepsilon_{\rm p}}{\varepsilon_{\rm h}},
\qquad
\mu_{\rm rel}\equiv\frac{\mu_{\rm p}}{\mu_{\rm h}},
\qquad
k_{\rm h}=\frac{\omega}{c}
\sqrt{\varepsilon_{\rm h}\mu_{\rm h}},
\qquad
m_{\rm r}=\frac{k_{\rm p}}{k_{\rm h}}
=\sqrt{\varepsilon_{\rm rel}\mu_{\rm rel}}.
\label{eq:host-wave-number}
\end{equation}
对被动介质取满足 $\operatorname{Im}k_j\ge0$ 的分支。准静态电偶极边值问题只涉及 $\varepsilon$；取通常的非磁性材料 $\mu_{\rm p}=\mu_{\rm h}=1$ 时即恢复
$m_{\rm r}=\sqrt{\varepsilon_{\rm p}/\varepsilon_{\rm h}}$。

先从精确球 Bessel 函数的最低阶幂看为什么“小粒子首先只剩偶极”。当 $|x|\ll1$ 时
\begin{align}
j_\ell(x)&=\frac{x^\ell}{(2\ell+1)!!}\left[1+O(x^2)\right],\\
y_\ell(x)&=-\frac{(2\ell-1)!!}{x^{\ell+1}}\left[1+O(x^2)\right].
\label{eq:smallx-spherical-bessel-general}
\end{align}
于是 Riccati--Bessel 函数满足
\begin{align}
\psi_\ell(x)=xj_\ell(x)&=O(x^{\ell+1}),\\
\xi_\ell(x)=xh_\ell^{(1)}(x)&=O(x^{-\ell})
\qquad (x\to0).
\end{align}
把这些阶数代回精确 Mie 分子和分母：对一般有限材料对比，分子含两个正则函数，最低为 $O(x^{2\ell+1})$；分母含一个正则函数和一个奇异外向函数，最低为 $O(1)$。因此一般地
\begin{equation}
\boxed{a_\ell,b_\ell=O(x^{2\ell+1})}
\qquad (\text{若相应低阶系数不因材料对称性额外消失}).
\label{eq:mie-multipole-scaling}
\end{equation}
这立刻说明 $\ell=1$ 偶极为 $O(x^3)$，$\ell=2$ 四极为 $O(x^5)$，依次至少再小两个 $x$ 的幂。因此 Rayleigh 近似不是“假定只有偶极”，而是精确部分波级数在 $x\to0$ 时的主导阶截断。

Rayleigh/quasistatic 条件首先是
\begin{equation}
\boxed{|k_{\rm h}|a\ll1,}
\label{eq:rayleigh-size-condition}
\end{equation}
其中 $a$ 为粒子尺度；若还要把电子附近的场直接用静电偶极场近似，则电子主要采样区域还应满足 $|k_{\rm h}|r\ll1$。

对任意形状的亚波长粒子，最低电偶极响应可以统一写成极化率张量
\begin{equation}
\bm p=\boldsymbol\alpha_0(\omega)\cdot\bm E_{\rm inc}(\bm r_0).
\label{eq:general-quasistatic-polarizability-def}
\end{equation}
在均匀宿主的准静态极限，偶极场由
\begin{equation}
\mathbf G_{\rm es}(\bm r,\bm r_0)
=\frac{1}{4\pi\varepsilon_0\varepsilon_{\rm h}}
\frac{3\hat{\bm R}\hat{\bm R}-\mathbf I}{R^3},
\qquad \bm R=\bm r-\bm r_0
\end{equation}
给出，所以同一个 PINEM 轨迹投影立即成为
\begin{equation}
\boxed{
\Ftilde
=\int\dd s\,\ee^{-\ii\kappa s}
\hat{\bm v}_0\cdot
\mathbf G_{\rm es}(\bm r_{\rm e}(s),\bm r_0)
\cdot\boldsymbol\alpha_0\cdot\bm E_{\rm inc}(\bm r_0).}
\label{eq:rayleigh-general-tensor-trajectory}
\end{equation}
因此 Rayleigh 层级真正新增的只是一个低维材料--几何对象 $\boldsymbol\alpha_0$；电子侧完全不变。任意椭球的退极化张量、$L_1+L_2+L_3=1$ 的证明以及动态辐射修正集中放在附录~\ref{app:ellipsoid-rayleigh}。正文下面直接回到与前一章精确 Mie 可以逐项核对的球。


\section{球：从静态极化率到偶极 Mie 渐近}
半径为 $a$ 的球置于均匀外场 $\bm E_{\rm inc}=E_{\rm L}\hat{\bm x}$。令 $\theta_x$ 为相对极化轴 $x$ 的极角。内外静电势的 $\ell=1$ 解只能写成
\begin{align}
\Phi_{\rm out}
&=-E_{\rm L}r\cos\theta_x+\frac{A\cos\theta_x}{r^2},\\
\Phi_{\rm in}
&=-B E_{\rm L}r\cos\theta_x.
\end{align}
$r=a$ 处切向电场连续等价于电势连续：
\begin{equation}
-E_{\rm L}a+\frac{A}{a^2}=-BE_{\rm L}a.
\label{eq:sphere-bc-general-1}
\end{equation}
法向位移连续给出
\begin{equation}
\varepsilon_{\rm h}
\left(-E_{\rm L}-\frac{2A}{a^3}\right)
=-\varepsilon_{\rm p}BE_{\rm L}.
\label{eq:sphere-bc-general-2}
\end{equation}
由第一式
\(
A=E_{\rm L}a^3(1-B)
\)，代入第二式：
\begin{align}
\varepsilon_{\rm h}\bigl[-1-2(1-B)\bigr]
&=-\varepsilon_{\rm p}B,\\
-3\varepsilon_{\rm h}+2\varepsilon_{\rm h}B
&=-\varepsilon_{\rm p}B,
\end{align}
故
\begin{equation}
B=\frac{3\varepsilon_{\rm h}}{\varepsilon_{\rm p}+2\varepsilon_{\rm h}},
\qquad
A=E_{\rm L}a^3
\frac{\varepsilon_{\rm p}-\varepsilon_{\rm h}}
{\varepsilon_{\rm p}+2\varepsilon_{\rm h}}.
\end{equation}
定义球的无量纲介电对比因子
\begin{equation}
\chi_{\rm sph}(\omega)
\equiv
\frac{\varepsilon_{\rm p}-\varepsilon_{\rm h}}
{\varepsilon_{\rm p}+2\varepsilon_{\rm h}},
\label{eq:chi-sphere-general}
\end{equation}
则偶极矩为
\begin{equation}
\boxed{
\bm p
=4\pi\varepsilon_0\varepsilon_{\rm h}a^3
\chi_{\rm sph}\,\bm E_{\rm inc}.}
\label{eq:sphere-polarizability-general}
\end{equation}
真空宿主 $\varepsilon_{\rm h}=1$ 时立即退化为熟悉的 Clausius--Mossotti 球极化率。

为检验准静态边值问题与精确 Maxwell 解的一致性，从 Mie 系数直接作同一个 $x=k_{\rm h}a\ll1$ 展开。以下计算在这里把已经声明的长波条件具体施加到精确 Mie 系数上。

精确 Mie 系数的低阶展开还必须同时检查电偶极与磁偶极通道，不能只凭“粒子很小”预先删除其中一个。
令
$x\ll1$，但此处不再假定 $\mu_{\rm rel}=1$。对 $\ell=1$，最低阶展开为
\begin{align}
\psi_1(x)
&=\frac{x^2}{3}-\frac{x^4}{30}+O(x^6),
&
\psi_1'(x)
&=\frac{2x}{3}-\frac{2x^3}{15}+O(x^5),
\\
\xi_1(x)
&=-\frac{\ii}{x}-\frac{\ii x}{2}
+\frac{x^2}{3}+O(x^3),
&
\xi_1'(x)
&=\frac{\ii}{x^2}-\frac{\ii}{2}
+\frac{2x}{3}+O(x^2).
\end{align}
对 electric/TM 系数，分子最低阶为
\begin{align}
N_1^{(E)}
&=m_{\rm r}\psi_1(m_{\rm r}x)\psi_1'(x)
-\mu_{\rm rel}\psi_1(x)\psi_1'(m_{\rm r}x)
\\
&=m_{\rm r}
\left(\frac{m_{\rm r}^2x^2}{3}\right)
\left(\frac{2x}{3}\right)
-\mu_{\rm rel}
\left(\frac{x^2}{3}\right)
\left(\frac{2m_{\rm r}x}{3}\right)
+O(x^5)
\\
&=\frac{2m_{\rm r}\mu_{\rm rel}}{9}
(\varepsilon_{\rm rel}-1)x^3+O(x^5),
\end{align}
因为 $m_{\rm r}^2=\varepsilon_{\rm rel}\mu_{\rm rel}$。分母最低阶为
\begin{align}
D_1^{(E)}
&=m_{\rm r}\psi_1(m_{\rm r}x)\xi_1'(x)
-\mu_{\rm rel}\xi_1(x)\psi_1'(m_{\rm r}x)
\\
&=m_{\rm r}
\left(\frac{m_{\rm r}^2x^2}{3}\right)
\frac{\ii}{x^2}
-\mu_{\rm rel}
\left(-\frac{\ii}{x}\right)
\left(\frac{2m_{\rm r}x}{3}\right)
+O(x^2)
\\
&=\frac{\ii m_{\rm r}\mu_{\rm rel}}{3}
(\varepsilon_{\rm rel}+2)+O(x^2).
\end{align}
所以
\begin{equation}
\boxed{
a_1
=-\frac{2\ii}{3}
\frac{\varepsilon_{\rm rel}-1}
{\varepsilon_{\rm rel}+2}x^3
+O(x^5).}
\label{eq:mie-rayleigh-a-general}
\end{equation}
magnetic/TE 通道也必须逐项展开。由式~\eqref{eq:mie-b-standard}，其分子为
\begin{align}
N_1^{(M)}
&=\mu_{\rm rel}\psi_1(m_{\rm r}x)\psi_1'(x)
-m_{\rm r}\psi_1(x)\psi_1'(m_{\rm r}x)\\
&=\mu_{\rm rel}
\left(\frac{m_{\rm r}^2x^2}{3}\right)
\left(\frac{2x}{3}\right)
-m_{\rm r}
\left(\frac{x^2}{3}\right)
\left(\frac{2m_{\rm r}x}{3}\right)
+O(x^5)\\
&=\frac{2m_{\rm r}^2}{9}(\mu_{\rm rel}-1)x^3+O(x^5).
\end{align}
分母为
\begin{align}
D_1^{(M)}
&=\mu_{\rm rel}\psi_1(m_{\rm r}x)\xi_1'(x)
-m_{\rm r}\xi_1(x)\psi_1'(m_{\rm r}x)\\
&=\mu_{\rm rel}
\left(\frac{m_{\rm r}^2x^2}{3}\right)
\frac{\ii}{x^2}
-m_{\rm r}
\left(-\frac{\ii}{x}\right)
\left(\frac{2m_{\rm r}x}{3}\right)
+O(x^2)\\
&=\frac{\ii m_{\rm r}^2}{3}(\mu_{\rm rel}+2)+O(x^2).
\end{align}
于是
\begin{equation}
\boxed{
b_1
=-\frac{2\ii}{3}
\frac{\mu_{\rm rel}-1}
{\mu_{\rm rel}+2}x^3
+O(x^5).}
\label{eq:mie-rayleigh-b-general}
\end{equation}
这表明一般磁介质在 Rayleigh 阶就有独立磁偶极响应。只有在非磁性特例 $\mu_{\rm rel}=1$ 时，分子中的 $x^3$ 项严格相消，因此必须把 $\psi_1,\psi_1'$ 多保留两阶：
\begin{align}
\psi_1(z)&=\frac{z^2}{3}-\frac{z^4}{30}+O(z^6),\\
\psi_1'(z)&=\frac{2z}{3}-\frac{2z^3}{15}+O(z^5).
\end{align}
于是
\begin{align}
\psi_1(m_{\rm r}x)\psi_1'(x)
&=\frac{2m_{\rm r}^2}{9}x^3
-\frac{2m_{\rm r}^2+m_{\rm r}^4}{45}x^5+O(x^7),\\
m_{\rm r}\psi_1(x)\psi_1'(m_{\rm r}x)
&=\frac{2m_{\rm r}^2}{9}x^3
-\frac{m_{\rm r}^2+2m_{\rm r}^4}{45}x^5+O(x^7).
\end{align}
两式相减，$x^3$ 项完全抵消，留下
\begin{equation}
N_1^{(M)}\big|_{\mu_{\rm rel}=1}
=\frac{m_{\rm r}^2(m_{\rm r}^2-1)}{45}x^5+O(x^7).
\end{equation}
而分母只需最低阶
\begin{equation}
D_1^{(M)}\big|_{\mu_{\rm rel}=1}
=\ii m_{\rm r}^2+O(x^2).
\end{equation}
故
\begin{equation}
\boxed{
b_1
=-\frac{\ii}{45}
(m_{\rm r}^2-1)x^5+O(x^7),
\qquad
\mu_{\rm rel}=1.}
\label{eq:mie-rayleigh-b-nonmagnetic}
\end{equation}
此时式~\eqref{eq:mie-rayleigh-a-general} 又化为
\begin{equation}
a_1
=-\frac{2\ii}{3}
\frac{m_{\rm r}^2-1}{m_{\rm r}^2+2}x^3+O(x^5),
\end{equation}
其中材料因子正是式~\eqref{eq:chi-sphere-general} 的
$\chi_{\rm sph}$。因此 quasistatic 电偶极结果就是完整 Mie electric/TM $\ell=1$ 通道的最低阶。


在 quasistatic 区，宿主中的散射近场是
\begin{equation}
\bm E_{\rm sc}(\bm r)
=E_{\rm L}\chi_{\rm sph}a^3
\frac{3\hat{\bm r}(\hat{\bm r}\cdot\hat{\bm x})-\hat{\bm x}}{r^3}.
\label{eq:sphere-quasistatic-field}
\end{equation}
取电子沿 $z$ 传播，横向位置写成
\begin{equation}
x=b\cos\varphi,
\qquad
y=b\sin\varphi,
\qquad r^2=b^2+z^2.
\end{equation}
于是电子真正耦合的纵向场为
\begin{equation}
E_{\omega,z}^{(\rm sph)}(b,\varphi,z)
=3E_{\rm L}\chi_{\rm sph}a^3
\frac{b\cos\varphi\,z}{(b^2+z^2)^{5/2}}.
\label{eq:Ez-sphere-general}
\end{equation}
定义电子相位匹配波数
\begin{equation}
\kappa\equiv\frac{\omega}{v_0}>0.
\label{eq:kappa-electron}
\end{equation}
利用
\begin{equation}
\frac{3z}{(b^2+z^2)^{5/2}}
=-\frac{\dd}{\dd z}(b^2+z^2)^{-3/2}
\end{equation}
并分部积分，边界项在 $|z|\to\infty$ 消失：
\begin{align}
\Ftilde^{(\rm sph)}
&=E_{\rm L}\chi_{\rm sph}a^3b\cos\varphi
\int_{-\infty}^{+\infty}\!\dd z\,
\frac{3z}{(b^2+z^2)^{5/2}}\ee^{-\ii\kappa z}\\
&=-\ii\kappa E_{\rm L}\chi_{\rm sph}a^3b\cos\varphi
\int_{-\infty}^{+\infty}\!\dd z\,
\frac{\ee^{-\ii\kappa z}}{(b^2+z^2)^{3/2}}.
\end{align}
剩余 Fourier 积分可以用 Schwinger 参数化逐步算出。对 $b>0$，有
\begin{equation}
\frac{1}{(b^2+z^2)^{3/2}}
=\frac{1}{\Gamma(3/2)}
\int_0^{\infty} t^{1/2}\ee^{-t(b^2+z^2)}\dd t
=\frac{2}{\sqrt\pi}
\int_0^{\infty} t^{1/2}\ee^{-t(b^2+z^2)}\dd t.
\label{eq:schwinger-three-halves}
\end{equation}
因为左端关于 $z$ 绝对可积，交换 $z$、$t$ 积分是合法的。于是
\begin{align}
I_{3/2}(b,\kappa)
&\equiv
\int_{-\infty}^{+\infty}
\frac{\ee^{-\ii\kappa z}}{(b^2+z^2)^{3/2}}\dd z\\
&=\frac{2}{\sqrt\pi}
\int_0^{\infty}t^{1/2}\ee^{-b^2t}
\left[
\int_{-\infty}^{+\infty}
\ee^{-tz^2-\ii\kappa z}\dd z
\right]\dd t.
\end{align}
内层是标准 Gaussian Fourier 积分。配方
\begin{equation}
-tz^2-\ii\kappa z
=-t\left(z+\frac{\ii\kappa}{2t}\right)^2
-\frac{\kappa^2}{4t},
\end{equation}
故
\begin{equation}
\int_{-\infty}^{+\infty}
\ee^{-tz^2-\ii\kappa z}\dd z
=\sqrt{\frac{\pi}{t}}
\exp\!\left(-\frac{\kappa^2}{4t}\right).
\end{equation}
因此
\begin{align}
I_{3/2}
&=2\int_0^{\infty}
\exp\!\left(-b^2t-\frac{\kappa^2}{4t}\right)\dd t.
\end{align}
再令 $u=b^2t$，则
\begin{align}
I_{3/2}
&=\frac{2}{b^2}
\int_0^{\infty}
\exp\!\left[-u-\frac{(\kappa b)^2}{4u}\right]\dd u.
\end{align}
使用修正 Bessel 函数的积分表示
\begin{equation}
K_\nu(x)
=\frac12\left(\frac{x}{2}\right)^\nu
\int_0^{\infty}u^{-\nu-1}
\exp\!\left(-u-\frac{x^2}{4u}\right)\dd u,
\end{equation}
并取 $\nu=-1$，再用 $K_{-1}=K_1$，得到
\begin{equation}
\int_0^{\infty}
\exp\!\left(-u-\frac{x^2}{4u}\right)\dd u
=xK_1(x).
\end{equation}
所以
\begin{equation}
\boxed{
\int_{-\infty}^{+\infty}
\frac{\ee^{-\ii\kappa z}}{(b^2+z^2)^{3/2}}\dd z
=\frac{2\kappa}{b}K_1(\kappa b).}
\label{eq:K1-trajectory-transform}
\end{equation}
代回上式得到
\begin{equation}
\boxed{
\Ftilde^{(\rm sph)}(b,\varphi)
=-2\ii E_{\rm L}\chi_{\rm sph}a^3
\kappa^2K_1(\kappa b)\cos\varphi.}
\label{eq:F-sphere-general}
\end{equation}
当 $\kappa b\ll1$ 时 $K_1(x)=x^{-1}+O(x\ln x)$，所以
\begin{equation}
|\Ftilde^{(\rm sph)}|
\simeq
2|E_{\rm L}\chi_{\rm sph}|a^3
\frac{\kappa}{b}|\cos\varphi|.
\label{eq:F-sphere-small-kappa}
\end{equation}
若实验主要采样 $b\sim a$，则 $|\Ftilde|\propto a^2$。


\section{材料、尺寸与电子相位匹配：三个彼此独立的控制量}

在均匀宿主 $\varepsilon_{\rm h}(\omega)$ 中，第一个独立控制量来自材料本身。对于本章正文保留的球几何，它就是式~\eqref{eq:chi-sphere-general} 已经由准静态边值问题和 Mie 渐近共同导出的
\begin{equation}
\boxed{
\chi_{\rm sph}(\omega)
=\frac{\varepsilon_{\rm p}(\omega)-\varepsilon_{\rm h}(\omega)}
{\varepsilon_{\rm p}(\omega)+2\varepsilon_{\rm h}(\omega)}.}
\label{eq:material-contrast-factors}
\end{equation}
因此准静态电偶极增强首先受分母
$\varepsilon_{\rm p}+2\varepsilon_{\rm h}$ 控制。若损耗有限而且实部可以穿过零点，则共振附近满足
\begin{equation}
\operatorname{Re}\!\left[\varepsilon_{\rm p}(\omega)+2\varepsilon_{\rm h}(\omega)\right]\simeq0,
\label{eq:sphere-quasistatic-resonance}
\end{equation}
同时还必须检查
$|\operatorname{Im}(\varepsilon_{\rm p}+2\varepsilon_{\rm h})|$
是否足够小；否则“实部为零”并不意味着大的可观测近场。只有在近似无损、$\varepsilon_{\rm h}$ 为实数时，式~\eqref{eq:sphere-quasistatic-resonance} 才进一步化成
$\operatorname{Re}\varepsilon_{\rm p}\simeq-2\varepsilon_{\rm h}$。

第二个独立控制量来自几何尺寸。精确 Mie 方程表明，决定长波展开是否成立的并不是粒子的绝对半径，而是无量纲 Maxwell 尺寸参数
\begin{equation}
\boxed{x\equiv k_{\rm h}a
=\frac{\omega a}{c}\sqrt{\varepsilon_{\rm h}(\omega)}.}
\label{eq:size-parameter-general}
\end{equation}
前文式~\eqref{eq:mie-rayleigh-a-general}--\eqref{eq:mie-rayleigh-b-general} 已经显示：$|x|\ll1$ 时电偶极 $\ell=1$ 是最低非零阶，而 $|x|\sim1$ 后高阶 $\ell=2,3,\ldots$ 不再受小参数压低，必须回到第~\ref{chap:mie-sphere}~章的完整 Mie 多极级数。此时多极共振与相位干涉是完整 Maxwell 边值问题的结果，不能用一个“修正过的静态极化率”代替。圆柱、椭球等其他几何的对应长波极限放在附录中；它们遵循同一逻辑，但不参与本章球几何的主推导。

第三个参数来自电子而不是 Maxwell 方程，即式~\eqref{eq:trajectory-tc} 中的相位匹配波数 $\kappa=\omega/v_0$。它决定轨迹泛函 $\mathcal L_{\rm e}$ 对空间 Fourier 分量的选择。故波长扫描同时改变
\begin{equation}
\varepsilon_{\rm p}(\omega),\qquad
\varepsilon_{\rm h}(\omega),\qquad
x=k_{\rm h}a,\qquad
\kappa=\omega/v_0,
\end{equation}
但四者物理来源不同。一个峰值可以来自材料极化共振、多极 Mie 共振，也可以来自散射近场的空间谱恰好在 $k_{\parallel}=\kappa$ 处增强；分析数据时不应把三者统称为同一种“共振”。偏振则通过入射场在相应多极通道中的展开系数以及轨迹投影 $\hat{\bm v}_0\cdot\bm E_\omega$ 共同进入。

